\chapter{PHÁT BIỂU BÀI TOÁN}\label{chap:baitoan}

\section{Bài toán}

Cho một kho bài báo khoa học tăng dần theo thời gian, cần xây dựng hệ thống cho phép: (i) tìm đúng bài cần đọc
bằng câu hỏi tự nhiên; (ii) hỏi đáp về nội dung một bài cụ thể; (iii) khám phá các bài liên quan. Mọi câu trả lời
phải có dẫn nguồn kiểm chứng được.

arXiv nhận hơn 20.000 bài mới mỗi tháng. Tìm kiếm theo từ khóa đơn thuần bỏ sót các bài dùng thuật ngữ khác; còn
chatbot dùng LLM thuần túy dễ bịa nội dung và bịa trích dẫn. Vì vậy hệ thống cần kết hợp ba thành phần: truy xuất
theo ngữ nghĩa, quan hệ giữa các bài (tác giả, chủ đề, trích dẫn), và LLM để tổng hợp câu trả lời từ đúng các
đoạn văn tìm được (RAG~\cite{rag}).

\section{Tác nhân và tình huống sử dụng}

\begin{table}[H]
\centering\small\setstretch{1.15}
\begin{tabularx}{\linewidth}{@{}lP{0.17\linewidth}LL@{}}
\toprule
\textbf{Mã} & \textbf{Tác nhân} & \textbf{Tình huống} & \textbf{Kết quả mong đợi}\\
\midrule
UC1 & Học viên, nhà nghiên cứu & Tìm bài báo bằng câu hỏi tự nhiên, ví dụ ``retrieval augmented generation cho văn bản pháp luật'' & Danh sách bài xếp hạng, có tiêu đề, tác giả, năm, link arXiv\\
UC2 & Nhà nghiên cứu & Tải lên một file PDF bài báo & Hệ thống báo ``đã có'' hoặc tự xử lý và thêm vào kho\\
UC3 & Nhà nghiên cứu & Tải lên hàng loạt PDF & Mỗi file có trạng thái riêng: thêm mới, trùng hoặc lỗi\\
UC4 & Nhà nghiên cứu & Hỏi về một bài cụ thể (phương pháp, dữ liệu, kết quả) & Câu trả lời dựa trên các đoạn của chính bài đó, có trích dẫn [n]\\
UC5 & Nhà nghiên cứu & Hỏi ``các bài liên quan đến bài này'' & Danh sách bài láng giềng trên đồ thị kèm lý do liên quan\\
UC6 & Người dùng & Trò chuyện thông thường, hỏi khái niệm chung & Chế độ Chat trả lời nhanh, không truy xuất\\
UC7 & Quản trị viên & Thu thập dữ liệu theo keyword hoặc chuyên mục, embedding trên GPU & Kho dữ liệu được mở rộng\\
\bottomrule
\end{tabularx}
\caption{Tác nhân và tình huống sử dụng}\label{tab:usecase}
\end{table}

\section{Phạm vi}

Hệ thống tập trung vào bài báo khoa học trên arXiv, viết bằng tiếng Anh, ở dạng PDF hoặc metadata. Câu hỏi có thể
bằng tiếng Việt hoặc tiếng Anh. Hệ thống không xử lý ảnh, bảng hay công thức trong PDF như dữ liệu có cấu trúc;
chúng chỉ được đọc như văn bản thường.
